% !TeX root = ../main.tex
%======================================================================
\chapter{Grundlagen der Lastflussberechnung}
\label{ch:grundlagen}
%======================================================================

Zur Berechnung von stationären Spannungen und Strömen in Stromnetzen und der Analyse von Leistungen ist die Lastflussberechnung in der Energietechnik eine der Grundlegenden Verfahren. In diesem Kapitel werden folgend die mathematischen Grundlagen eingeführt, die als Basis für diese Ausarbeitung verwendet werden.


%Die Lastflussberechnung (engl. \textit{Power Flow}, PF) ist eine der fundamentalsten Analysen in der elektrischen Energietechnik. Sie dient der Bestimmung der stationären Spannungen und Ströme an allen Knoten eines Netzes unter gegebenen Erzeugungs- und Verbrauchsbedingungen. Dieses Kapitel führt die mathematischen Grundlagen ein, auf denen die in \cref{ch:tpf,ch:methoden_pv} vorgestellten Methoden aufbauen.

\section{Netzmodell und Knotentypen}
\label{sec:netzmodell}

Bevor die Berechnung von Lastflüssen eingeführt wird, ist zunächst eine Modellierung von elektrischen Netzen notwendig.
Diese folgt der Notation in \cite{Salazar2024} und \cite{Giraldo2022}, da sich die Tensorformulierung, als auch Implementierung in Kapitel \cref{ch:tpf,ch:methoden_pv} nach dieser richtet.  

%Dieses Kapitel führt die grundlegende Modellierung eines elektrischen Verteilnetzes für die Lastflussberechnung ein. Alle in den folgenden Kapiteln verwendeten Symbole und Konventionen werden hier definiert. Die Darstellung orientiert sich an der in \cite{Salazar2024} und \cite{Giraldo2022} verwendeten Notation, die als Grundlage für die spätere Tensorformulierung dient.

\subsection{Admittanzformulierung des Netzes}
\label{subsec:admittanz}

Ein elektrisches Netz kann als Graph mit einer Gesamtknotenmenge $\Omega_\mathrm{B}$ modelliert werden. Diese besteht aus Quellknoten (Slack-Knoten), die eine Energiezufuhr modellieren und in der Menge $\Omega_\mathrm{s}$ enthalten sind. Die restlichen Knoten sind Lastknoten (\emph{drain}) und haben die Menge $\Omega_\mathrm{d}$. Diese Mengen sind disjunkt, sodass $\Omega_B = \Omega_s \cup \Omega_d$ und $\Omega_s \cap \Omega_d = \emptyset$. Zur kompakteren Schreibweise beschreibt $b = |\Omega_\mathrm{d}|$ die Anzahl der Lastknoten.\\

An allen Knoten gelten das Ohm'sche und Kirchoff'sche Gesetz welches erweitert auf ein Multiknotensystem durch das Vektor-Matrix-System

%Ein elektrisches Netz wird als gerichteter Graph mit einer Knotenmenge $\Omega_B$ modelliert. Diese wird disjunkt zerlegt in die Menge der Slack-Knoten $\Omega_s$ und die Menge der Lastknoten $\Omega_d$, sodass $\Omega_B = \Omega_s \cup \Omega_d$ und $\Omega_s \cap \Omega_d = \emptyset$ gilt. Sei $b = |\Omega_d|$ die Anzahl der Lastknoten.
%
%Die Beziehung zwischen den komplexen Knotenspannungen $\mathbf{v} \in \mathbb{C}^{|\Omega_B|}$ und den komplexen Knotenstromeinspeisungen $\mathbf{i} \in \mathbb{C}^{|\Omega_B|}$ wird durch die Knotenadmittanzmatrix $\mathbf{Y} \in \mathbb{C}^{|\Omega_B| \times |\Omega_B|}$ nach dem Kirchhoffschen Knotensatz beschrieben. Durch Umsortierung der Knoten in Slack- und Lastknoten kann die Admittanzmatrix in Blockform dargestellt werden~\cite{Salazar2024}:
\begin{equation}
	\label{eq:admittance_block}
	\begin{bmatrix} \mathbf{i}_s \\ -\mathbf{i}_d \end{bmatrix}
	=
	\begin{bmatrix} \mathbf{Y}_{ss} & \mathbf{Y}_{sd} \\ \mathbf{Y}_{ds} & \mathbf{Y}_{dd} \end{bmatrix}
	\begin{bmatrix} \mathbf{v}_s \\ \mathbf{v}_d \end{bmatrix}
\end{equation}

beschrieben wird. 
Dabei stehen komplexen Knotenspannungen $\mathbf{v} \in \mathbb{C}^{|\Omega_\mathrm{B}|}$ und $\mathbf{i} \in \mathbb{C}^{|\Omega_\mathrm{B}|}$ über die Admittanzmatrix $\mathbf{Y} \in \mathbb{C}^{|\Omega_\mathrm{B}| \times |\Omega_\mathrm{B}|}$ in Beziehung \cite{TinneyHart1967, Stott1974, Bernstein2018, Chen1991}. \todo{literatur} 
Aufgespalten nach den eingeführten Knotenmengen bezeichnet $\mathbf{v}_\mathrm{s} \in \mathbb{C}^{|\Omega_\mathrm{s}|}$ die bekannten Quellspannungen und $\mathbf{v}_\mathrm{d} \in \mathbb{C}^{b}$ die unbekannten Spannungen an den Lastknoten. Das negative Vorzeichen bei $\mathbf{i}_d$ folgt aus der Konvention, dass Quellen positiv, jedoch Lasten negativ gewertet werden \cite{Kersting2017, Milano2010}. \\ \todo{literatur}

Aus Gleichung \eqref{eq:admittance_block} lässt sich der Ausdruck für den Laststrom 

%Aus dem unteren Block von \eqref{eq:admittance_block} folgt die zentrale Gleichung für die Lastknoten:
\begin{equation}
	\label{eq:current_balance}
	-\mathbf{i}_d = \mathbf{Y}_{ds}\mathbf{v}_s + \mathbf{Y}_{dd}\mathbf{v}_d
\end{equation}

für die nachfolgende Fixpunktiteration isolieren.
Die komplexe Knotenleistung $\mathbf{s}_\mathrm{d}$ und Knotenstrom $\mathbf{i}_\mathrm{d}$ ist elementweise durch
\begin{equation}
	\label{eq:power_current}
	\mathbf{s}_d = \mathbf{v}_d \odot \mathbf{i}_d^{*}
	\quad\Longleftrightarrow\quad
	\mathbf{i}_d = \left(\mathbf{s}_d \oslash \mathbf{v}_d\right)^{*}
\end{equation}
charakterisiert\todo{literatur}, wobei $\odot$ das Hadamard-Produkt und $\oslash$ die Hadamard-Division und $(\cdot)^{*}$ die komplexe Konjugation darstellt~\cite{Grainger1994, Milano2010, Bernstein2018}. 
Die Nichtlinearität des Lastflussproblems entsteht durch das Produkt des Stroms mit der Spannung in \eqref{eq:power_current}. 
Eingesetzt in \eqref{eq:current_balance} entsteht die Grundgleichung des Lastflusses
%Einsetzen in \eqref{eq:current_balance} liefert die Grundgleichung des Lastflusses:
\begin{equation}
	\label{eq:pf_fundamental}
	-\operatorname{diag}(\mathbf{v}_d^{*})^{-1}\mathbf{s}_d^{*} = \mathbf{Y}_{ds}\mathbf{v}_s + \mathbf{Y}_{dd}\mathbf{v}_d
\end{equation}

welche allgemein ist und gilt sowohl für beliebig-phasige als auch für radiale und vermaschte Netze~\cite{Salazar2024}. Sie bildet den Ausgangspunkt für die in Kapitel~\ref{ch:tpf} vorgestellte Fixpunktiteration.

\subsection{Knotentypen}
\label{subsec:knotentypen}

Jeder elektrische Knoten verfügt über vier Größen: der Wirkleistung $P$, der Blindleistung $Q$, dem Spannungsbetrag $|V|$ und dem Spannungswinkel $\theta$. 
Davon sind immer zwei für den Lastfluss bekannt und die Klassifikation von Lastknoten erfolgt über deren gesuchten und gegebenen Größen. 
Auf dieser Basis folgen die in Tabelle~\ref{tab:knotentypen} vorgestellten drei Knotentypen: PQ-, PV- und Slack-Knoten.\\ 

%In der Lastflussberechnung werden Netzknoten anhand der jeweils vorgegebenen und gesuchten Größen klassifiziert. Jeder Knoten besitzt vier elektrische Größen -- Wirkleistung $P$, Blindleistung $Q$, Spannungsbetrag $|V|$ und Spannungswinkel $\theta$ --, von denen im Lastfluss zwei vorgegeben und zwei unbekannt sind. Daraus ergeben sich drei Standardtypen, dargestellt in Tabelle~\ref{tab:knotentypen}.\\

\begin{table}[h]
	\centering
	\caption{Knotentypen im Lastflussproblem unter gesuchten und gegebenen Größen mit Anwendungsbeispielen~\cite{TinneyHart1967, Stott1974, Milano2010, Zimmerman2011}}
	\label{tab:knotentypen}
	\begin{tabular}{lccp{8.9cm}}
		\toprule
		Typ & Bekannt & Unbekannt & Typische Netzelemente \\
        \hline\hline
		Slack & $|V|,\theta$ & $P,Q$ &{ \footnotesize Umspannwerk zum überlagerten Netz (Netzverknüpfungspunkt), Referenzgenerator im Übertragungsnetz }\\
		\midrule
		PQ    & $P,Q$        & $|V|,\theta$ & {\footnotesize Haushalts- und Gewerbelasten, Industrieabnehmer, kleine PV-Anlagen im $\cos\phi$-Betrieb, ungeregelte Windkraftanlagen, konventionelle Verbraucher} \\
		\midrule
		PV    & $P,|V|$      & $Q,\theta$ & {\footnotesize Synchrongeneratoren mit Erregerregelung (Wasser-, Gas-, Dampfkraftwerke), größere PV-Anlagen und Batteriewechselrichter mit Q(U)-Regelung nach VDE-AR-N 4110/4120, STATCOM-Anlagen} \\
		\bottomrule
	\end{tabular}
\end{table}
\todo{literatur.}

%Der Slack-Knoten dient als Winkel- und Spannungsreferenz des Netzes. Sein Spannungsbetrag und -winkel sind fest vorgegeben (typischerweise $|V|=1{,}0$~p.u. und $\theta = 0$). Er bilanziert die im Netz auftretenden Verluste sowie die Differenz zwischen Erzeugung und Verbrauch, weshalb seine Wirk- und Blindleistungseinspeisungen erst nach Lösung des Lastflusses berechnet werden können. In Verteilnetzen entspricht der Slack-Knoten dem Netzverknüpfungspunkt zum überlagerten Netz.\\
Slack-Knoten sind Einspeisungsknoten und fungieren als Spannungs- und Winkelreferenz eines Netzes. 
Hier sind dementsprechend Spannungswinkel und -betrag konstruktiv vorgegeben und der Slack-Knoten bilanziert die Differenz zwischen Erzeugung und Verbrauch.
Dadurch ist es weiterhin erst nach der Lastflussberechnung möglich die Einspeisung von Wirk- und Blindleistungen zu bestimmen. 
In den hier verwendeten Verteilernetzen modelliert der Slack-Knoten die Anbindung zur übergeordneten Netzebene~\cite{Stott1974, Milano2010, Kundur1994}.\\
\todo{literatur.}

%An PQ-Knoten sind Wirk- und Blindleistung vorgegeben, während Spannungsbetrag und -winkel gesucht sind. Sie modellieren typische Lasten sowie Einspeiser ohne Spannungsregelung (z.\,B.\ kleine PV-Anlagen im $\cos\varphi$-Betrieb). Die überwiegende Mehrheit der Knoten in Verteilnetzen ist von diesem Typ.\\
PQ-Knoten modellieren typische Lasten in einem elektrischen Netz, wie z.B. kleine PV-Anlagen im $\cos\phi$-Betrieb, oder Haushaltsgeräte. An solchen Knoten ist die Wirk- und Blindleistung vorgegeben und es wird der Spannungsbetrag und -winkel im Lastfluss berechnet. Dieser Knotentyp ist für gewöhnlich am meißten vertreten~\cite{Kersting2017, Milanovic2013}.\\
\todo{literatur.}

%An PV-Knoten sind Wirkleistung und Spannungsbetrag vorgegeben, während Blindleistung und Spannungswinkel unbekannt sind. Sie modellieren Synchrongeneratoren mit Erregerregelung sowie größere PV-Anlagen und Batteriewechselrichter mit Spannungshaltungsfunktion. Physikalisch stellt der PV-Knoten eine Blindleistungsquelle bzw. -senke dar, die soviel $Q$ einspeist oder aufnimmt, wie zur Einhaltung des Spannungssollwertes $|V_k^{\text{spec}}|$ nötig ist.\\
PV-Knoten modellieren Synchrongeneratoren mit Erregerregelung sowie größere PV-Anlagen und Batteriewechselrichter mit Spannungshaltungsfunktion. Dort sind der Spannungsbetrag und Wirkleistung vorgegeben. Dadurch ist ein PV-Knoten eine Blindleistungsquelle/-senke, da sie genau so viel Blindleistung erzeugt oder aufnimmt, um die nötige Spannung $|V_k^{\text{spec}}|$ und Wirkleistung aufrecht zu erhalten~\cite{daCosta1999, Kundur1994, IEEE1547_2018}.\\
\todo{literatur.}

%In der Praxis ist die Blindleistungseinspeisung durch die Kapazität der Regeleinrichtung begrenzt:
%\begin{equation}
%	\label{eq:qlimits}
%	Q_k^{\min} \leq Q_k \leq Q_k^{\max}
%\end{equation}
%Wird eine Grenze erreicht, so wird der Knoten üblicherweise in einen PQ-Knoten mit $Q_k = Q_k^{\min}$ bzw.\ $Q_k = Q_k^{\max}$ umgewandelt und die Spannungsbedingung $|V_k|=|V_k^{\text{spec}}|$ wird dann fallengelassen. Diese Umschaltung wird in Kapitel~\ref{sec:pv_methoden} bei den vorgeschlagenen Methoden berücksichtigt.\\ \todo{literatur.}

%Die Behandlung der PV-Knoten ist der methodische Kern dieser Arbeit: Während sie sich in das Newton-Raphson-Verfahren (Abschnitt~\ref{sec:nr}) auf natürliche Weise über eine Modifikation der Jacobi-Matrix einfügen lassen~\cite{daCosta1999}, bricht bei ihrer Übertragung auf die Fixpunkt-Iteration die zentrale Voraussetzung einer konstanten Systemmatrix, zusammen. Dies motiviert die in dieser Arbeit entwickelten Methode.

Diese Verteilung der gegebenen und gesuchten Größen motiviert den Kern dieser Ausarbeitung. 
PV-Knoten setzen andere Informationen über das Netz voraus und bedürfen einer eigenen Handhabung in den Lösungsverfahren. 
In Abschnitt~\ref{sec:nr} wird deren Behandlung im Newton-Raphson (NR) eingeführt.
Dort wird aufgezeigt wieso die Implementierung im NR über eine simple Modifikation in der Jakobi-Matrix erfolgen kann, wohingegen die Implementierung in der Fixpunktiteration zu dem zentralen Problem dieser Arbeit führt. 

\subsection{Lastmodellierung}
\label{subsec:lastmodell}

%Reale Lasten sind spannungsabhängig. Zur Modellierung dieser Abhängigkeit wird das ZIP-Modell verwendet, das die Last als Superposition eines konstanten Impedanzanteils ($Z$), eines konstanten Stromanteils ($I$) und eines konstanten Leistungsanteils ($P$) beschreibt~\cite{Giraldo2022}:
Elektrische Lasten sind weiterhin selber spannungsabhängig. 
Den Industriestandard bildet das ZIP-Modell welches eine Last durch eine Superposition eines konstanten Impedanzanteils ($Z$), eines konstanten Stromanteils ($I$) und eines konstanten Leistungsanteils ($P$) als eine Funktion der Spannung beschreibt~\cite{Giraldo2022}. \todo{literatur} 
Die Leistung setzt sich dementsprechend durch

\begin{equation}
	\label{eq:zip}
	\mathbf{s}_d = \left(\boldsymbol{\alpha}_P + \boldsymbol{\alpha}_I \odot |\mathbf{v}_d| + \boldsymbol{\alpha}_Z \odot |\mathbf{v}_d|^2\right) \odot \mathbf{s}_n
\end{equation}

zusammen. $\mathbf{s}_n$ ist hier die Nennleistung der Last und $\boldsymbol{\alpha}_P, \boldsymbol{\alpha}_I, \boldsymbol{\alpha}_Z$ sind die vektoriellen knotenweisen ZIP-Koeffizienten. 
Dabei gilt, dass die Summe der Koeffizienten eines Knotens sich zu eins aufaddieren 

\begin{equation}
	\alpha_{P,k} + \alpha_{I,k} + \alpha_{Z,k} = 1 \quad \forall k \in \Omega_d
\end{equation}

und das Modell damit bloß die Aufteilung der Spannungsabhängigkeit beschreibt.
Der Fall $\boldsymbol{\alpha}_P = \mathbf{1}$, $\boldsymbol{\alpha}_I = \boldsymbol{\alpha}_Z = \mathbf{0}$ entspricht einem konstanten Leistungsmodell, welches in dieser Arbeit als Standardfall angenommen wird.
Diese Annahme ist in stationären Netzen und in Optimierungsrechnungen gängig, da konstante Lasten keine Selbstregulierung aufweisen und somit die maximalen Spannungsabfälle und Verluste vorweisen und damit den ungünstigsten Fall darstellen \cite{Kundur1994,Kersting2017,IEEETF1995,VanCutsem1998,Bolognani2016}. \todo{check}



\subsection{Thévenin-Äquivalent an einem Lastknoten}
\label{subsec:thevenin}


\begin{figure}[htbp]
	\centering
	% Falls zu breit: \resizebox{\linewidth}{!}{ ... }
	\resizebox{5.91in}{!}{
	\begin{tikzpicture}[
		gen/.style={circle, draw=black, minimum size=5mm, font=\bfseries\small, line width=0.4pt},
		bus/.style={line width=3.5pt},
		loadarrow/.style={-{Triangle[length=2.5mm,width=2.5mm]}, line width=1pt},
		wire/.style={line width=0.9pt},
		>=Latex
		]
		
		%=========================================================
		% LEFT PANEL: "The rest of the power system"
		%=========================================================
		\begin{scope}
			\draw[dashed, line width=1pt, rounded corners=2pt]
			(-0.4,-0.3) rectangle (6.2,5.9);
			
			% ---- Bus bars (each a short thick horizontal segment) ----
			\draw[bus] (0.6,4.3) -- (1.8,4.3);   % Bus A
			\draw[bus] (3.2,4.6) -- (4.4,4.6);   % Bus B
			\draw[bus] (1.6,2.9) -- (3.4,2.9);   % Bus C (hub)
			\draw[bus] (0.3,1.3) -- (1.5,1.3);   % Bus D
			\draw[bus] (3.3,1.5) -- (4.9,1.5);   % Bus E
			
			% ---- Generators (stub up from bus to G) ----
			\node[gen] (g1) at (1.2,5.1) {G};
			\draw[wire] (g1) -- (1.2,4.3);
			\node[gen] (g2) at (3.8,5.4) {G};
			\draw[wire] (g2) -- (3.8,4.6);
			\node[gen] (g3) at (2.3,3.7) {G};
			\draw[wire] (g3) -- (2.3,2.9);
			\node[gen] (g4) at (4.1,2.3) {G};
			\draw[wire] (g4) -- (4.1,1.5);
			
			% ---- Loads (stub down from bus, arrow) ----
			\draw[wire, loadarrow] (0.9,1.3) -- (0.9,0.4);
			\draw[wire, loadarrow] (4.6,1.5) -- (4.6,0.6);
			
			% ---- Transmission lines connecting bus bars (mesh) ----
			\draw[wire] (1.8,4.3) -- (1.6,2.9);   % A -- C
			\draw[wire] (3.2,4.6) -- (2.8,2.9);   % B -- C
			\draw[wire] (2.0,2.9) -- (1.5,1.3);   % C -- D
			\draw[wire] (3.3,1.5) -- (1.5,1.3);   % D -- E  (bottom tie line)
			
			% ---- outgoing feeders toward the external "Load node k" bus ----
			\draw[wire] (3.4,2.9) -- (5.4,2.9);   % C's feeder, exits box at (6.2,*)
			\draw[wire] (4.9,1.5) -- (5.4,1.5);   % E's feeder, exits box at (6.2,*)
			\draw[dotted, line width=0.9pt] (5.4,2.9) -- (7.0,3.1);
			\draw[dotted, line width=0.9pt] (5.4,1.5) -- (7.0,2.3);
			
			% ---- caption ----
			\node[font=\bfseries] at (2.9,-0.7) {Restliches Netz};
			
		\end{scope}
		
		%=========================================================
		% "Load node k" bus, sitting between the two panels
		%=========================================================
		\draw[bus] (7.0,2.0) -- (7.0,3.3);
		\draw[wire, loadarrow] (7.0,2.0) -- (7.0,1.3);
		\node[above] at (7.0,3.75) {\small Last $k$};
		
		%=========================================================
		% GREEN TRANSFORMATION ARROW
		%=========================================================
		\draw[-{Triangle[length=4mm,width=5mm]}, line width=3pt]
		(7.5,2.35) -- (8.9,2.35);
		
		%=========================================================
		% RIGHT PANEL: Thevenin equivalent circuit
		%=========================================================
		\begin{scope}[shift={(9.3,0)}]
			
			\draw[dashed, line width=1pt, rounded corners=2pt]
			(-0.3,1.2) rectangle (4.6,3.5);
			
			% --- Thevenin voltage source (circle with sine wave) ---
			\node[circle, draw=black, minimum size=1.5cm, line width=0.8pt] (src) at (0.9,2.35) {};
			\draw[wire, line width=1pt] plot[smooth, tension=1] coordinates {
				($(src.center)+(-0.42,0)$)   ($(src.center)+(-0.21,0.28)$)
				($(src.center)+(0,0)$)       ($(src.center)+(0.21,-0.28)$)
				($(src.center)+(0.42,0)$) };
			\node[below=2pt] at (0.9,1.3) {$E_{th,k}$};
			
			% wire from source to impedance
			\draw[wire] (1.65,2.35) -- (2.1,2.35);
			
			% --- Thevenin impedance (rectangle) ---
			\draw[line width=0.8pt] (2.1,2.1) rectangle (3.5,2.6);
			\node[below=2pt] at (2.8,1.3) {$Z_{th,k}$};
			
			% wire from impedance to the boundary of the dashed box
			\draw[wire] (3.5,2.35) -- (4.6,2.35);
			
		\end{scope}
		% (dashed box right edge is now at global x = 9.3 + 4.6 = 13.9)
		
		% wire continuing from the dashed box out to the load bus, with current label
		\draw[wire] (13.9,2.35) -- (15.4,2.35);
		\node[above] at (14.3,2.5) {\small $I_{L,k}$};
		
		% --- Load node k bus on the far right, with load arrow and voltage label ---
		\draw[bus] (15.4,2.0) -- (15.4,2.7);
		\draw[wire, loadarrow] (15.4,2.0) -- (15.4,1.1);
		\node[above right, xshift=-2pt] at (15,2.75) {\small Last $k$};
		\node[right] at (15.55,1.55) {\small $U_{L,k}$};
		
	\end{tikzpicture}
	}
	\caption[Th\'evenin-\"Aquivalent an einem Lastknoten]{Reduktion des Netzes auf
		das Th\'evenin-\"Aquivalent aus Sicht des Lastknotens $k$. Das gesamte
		\"ubrige Netz einschlie{\ss}lich Slack-Einspeisung und aller weiteren Lasten
		wird durch die Leerlaufspannung $E_{\mathrm{Th},k}$ und die Innenimpedanz
		$Z_{\mathrm{Th},k}=[\mathbf{Z}_B]_{kk}=R_{kk}+\mathrm{j}X_{kk}$
		ersetzt~\cite{Grainger1994,Anderson1995,Dorfler2013}.}
	\label{fig:thevenin}
\end{figure}

Für die spätere Herleitung der Q-Sensitivität für die Äußere Schleife wird das Thévenin-Äquivalent an einem Knoten benötigt. Aus der Sicht eines Knotens $k \in \Omega_d$ lässt sich das restliche Netz gemäß Abbildung \ref{fig:thevenin} durch eine ideale Spannungsquelle $E_{\text{Th},k}$ in Reihe mit einer Impedanz $Z_{\text{Th},k}$ ersetzen~\cite{osipov2022}.
Die Bus-Impedanzmatrix ist definiert als
\begin{equation}
	\label{eq:zbus}
	\mathbf{Z}_B := \mathbf{Y}_{dd}^{-1}.
\end{equation}
Ihre Diagonalelemente entsprechen den Thévenin-Impedanzen der jeweiligen Knoten \cite{Giraldo2022,ALDAOUDEYEH2019105890}
\begin{equation}
	\label{eq:z_thev}
	Z_{\text{Th},k} = R_{kk} + jX_{kk} = [\mathbf{Z}_B]_{kk}
\end{equation}
Dieser Zusammenhang liefert die theoretische Grundlage für die in Methode A verwendete Sensitivitätsapproximation für einen gegebenen Arbeitspunkt.
%\begin{equation}
%	\label{eq:q_sensitivity}
%	\frac{\partial |V_k|^2}{\partial Q_k} \approx 2 X_{kk}
%\end{equation}
%welche in Kapitel~\ref{sec:pv_methoden} hergeleitet wird. In Verteilnetzen mit hohem $R/X$-Verhältnis ist zu beachten, dass $R_{kk}$ nicht vernachlässigbar ist, weshalb dort ggf.\ die vollständige komplexe Sensitivität verwendet werden muss.

%Aus \eqref{eq:current_balance} folgt mit \eqref{eq:zbus}
%\begin{equation}
%	\label{eq:v_zbus}
%	\mathbf{v}_d = \underbrace{-\mathbf{Z}_B\mathbf{Y}_{ds}\mathbf{v}_s}_{=:\ \mathbf{v}_d^{0}}
%	\;-\; \mathbf{Z}_B\mathbf{i}_d ,
%\end{equation}
%wobei $\mathbf{v}_d^{0}$ die Leerlaufspannungen der Lastknoten bezeichnet. Die
%$k$-te Zeile von \eqref{eq:v_zbus} lautet
%\begin{equation}
%	\label{eq:thev_derivation}
%	v_k = \underbrace{v_k^{0} - \sum_{j \in \Omega_d \setminus \{k\}} [\mathbf{Z}_B]_{kj}\, i_j}_{=:\ E_{\mathrm{Th},k}}
%	\;-\; \underbrace{[\mathbf{Z}_B]_{kk}}_{=:\ Z_{\mathrm{Th},k}}\, i_k ,
%\end{equation}
%was genau der Klemmengleichung einer Spannungsquelle $E_{\mathrm{Th},k}$ in Reihe
%mit $Z_{\mathrm{Th},k}$ entspricht~\cite{Grainger1994,Anderson1995,Dorfler2013}.

%\subsection{Existenz einer Lösung}
%\label{subsec:existenz}
%Die Konvergenz einer Lastflussrechnung ist nicht immer möglich. Die von Bolognani et. al.~\cite{bolognani2016} angeführte Analyse eines Zweiknotennetzes mit Slack-Spannung $v_0$, Quellimpedanz $z_s$ und Lastleistung $s_l$ ergibt, dass es physikalische Lösungen gibt, wenn
%%Nicht jeder Lastzustand führt zu einem konvergenten Lastfluss. Nach den Ergebnissen von Bolognani und Zampieri~\cite{bolognani2016} existiert für ein Zweiknotensystem mit Slack-Spannung $v_0$, Quellimpedanz $z_s$ und Lastleistung $s_l$ genau dann eine praktische (Hochspannungs-)Lösung, wenn
%\begin{equation}
%	\label{eq:feasibility}
%	\|v_0\|^2 \geq 4\,\|s_l\|\,\|z_s\|
%\end{equation}
%gilt. Der Grenzfall entspricht dem Punkt maximaler Leistungsübertragung, an dem die zwei Impedanzlösungen $z_{l,1}$ und $z_{l,2}$ zusammenfallen. Eine geometrische Interpretation dieser Bedingung sowie ihre Verallgemeinerung auf Mehrknotensysteme liegt in \cite{Salazar2024} vor.

%Diese Bedingung ist von Bedeutung, weil sie im Grenzfall zusammenfällt mit der Kontraktionsbedingung des Banachschen Fixpunktsatzes (Abschnitt~\ref{sec:fpi}): Nur unterhalb des MLP existiert eine physikalisch sinnvolle Lösung, und nur dort konvergiert die Fixpunkt-Iteration.


%\bigskip
%\noindent
%Zusammenfassend wurden in diesem Abschnitt die Netzmodellierung über die Blockformulierung der Admittanzmatrix, die drei Knotentypen mit ihren jeweiligen bekannten und unbekannten Größen, das ZIP-Lastmodell, das Thévenin-Äquivalent und die grundlegende Existenzbedingung für Lastflusslösungen eingeführt. Auf dieser Basis werden in den folgenden Abschnitten das Newton-Raphson-Verfahren (Abschnitt~\ref{sec:nr}) und die Fixpunkt-Iteration (Abschnitt~\ref{sec:fpi}) als konkrete Lösungsverfahren dargestellt.









\section{Newton-Raphson-Verfahren}
\label{sec:nr}

Ein weit verbreitetes Verfahren in der Lastflussberechnung bildet das numerische Newton-Raphson-Verfahren (NR-Verfahren)\todo{literatur}.
Die Basis hierfür bildet eine Leistungsbilanz an jedem einzelnen Lastknoten $k$. 
Das NR-Verfahren kann generell in kartesischen oder in polaren Koordinaten formuliert werden.
Zur kompakteren Übersicht wird das Verfahren hier in polaren Koordinaten eingeführt.
Für die Leistungsbilanz werden die Knotenspannungen dementsprechend durch

%Ausgangspunkt ist die Leistungsbilanz an jedem Lastknoten $k \in \Omega_d$. Im Gegensatz zur in Abschnitt~\ref{subsec:admittanz} verwendeten kartesischen Formulierung ($\mathbf{v} \in \mathbb{C}$) wird für das NR-Verfahren die polare Darstellung der komplexen Knotenspannungen bevorzugt, da sich die Leistungsflussgleichungen dadurch kompakter formulieren lassen. Jede Knotenspannung wird zerlegt in Betrag und Winkel:
\begin{equation}
	\label{eq:polar_voltage}
	V_k = |V_k| \, e^{j\theta_k}, \qquad k \in \Omega_B
\end{equation}

dargestellt, $|V_k| \in \mathbb{R}_{+}$ den Spannungsbetrag und $\theta_k \in [-\pi, \pi]$ den Spannungswinkel am Knoten $k$ bezeichnet.\todo{literatur.}
Für Mehrknotensysteme werden die Größen vektoriell zu $|\mathbf{V}|$ und $\boldsymbol{\theta}$ für alle Netzknoten zusammengefasst.
%Die Vektoren $|\mathbf{V}|$ und $\boldsymbol{\theta}$ fassen die Spannungsbeträge bzw. -winkel aller Knoten zusammen.
Die komplexe Knotenadmittanzmatrix $\mathbf{Y}$ \eqref{eq:admittance_block} in ihrer Real- und Imaginärzerlegung 
%Analog wird die komplexe Knotenadmittanzmatrix $\mathbf{Y}$ aus \eqref{eq:admittance_block} in Real- und Imaginärteil zerlegt:
\begin{equation}
	\label{eq:y_split}
	\mathbf{Y} = \mathbf{G} + j\mathbf{B}, \qquad Y_{km} = G_{km} + jB_{km}
\end{equation}

führt die Konduktanzmatrix $\mathbf{G} \in \mathbb{R}^{|\Omega_B|\times|\Omega_B|}$ und die Suszeptanzmatrix $\mathbf{B} \in \mathbb{R}^{|\Omega_B|\times|\Omega_B|}$ ein.
%mit der Konduktanzmatrix $\mathbf{G} \in \mathbb{R}^{|\Omega_B|\times|\Omega_B|}$ und der Suszeptanzmatrix $\mathbf{B} \in \mathbb{R}^{|\Omega_B|\times|\Omega_B|}$. 
Die Diagonalelemente bilden damit die Selbstadmittanz eines Knotens und bestehen aus allen an diesem angeschlossenen Einzeladmittanzen \todo{literatur.}. 
Folglich bilden Nebendiagonalelemente die elektrischen Kopplungen zwischen den Knoten $k$ und $m$.
Die Lei,stungsbilanzen selber werden durch die in den PQ-Knoten spezifizierten Wirk- und Blindleistungsvorgaben $P_k^{\text{spec}}$, $Q_k^{\text{spec}}$ und den berechneten Leistungswerten $P_k^{\text{calc}}$ und $Q_k^{\text{calc}}$ gebildet.
Deren Differenz
%Das Element $Y_{km}$ beschreibt die Kopplung zwischen den Knoten $k$ und $m$. 
%Für jeden Lastknoten $k \in \Omega_d$ sind die Wirk- und Blindleistungsvorgaben $P_k^{\text{spec}}$ bzw.\ $Q_k^{\text{spec}}$ gegeben (bei PQ-Knoten aus dem Lastprofil, bei PV-Knoten nur $P_k^{\text{spec}}$. Aus \eqref{eq:pf_fundamental} folgen die Leistungsungleichungen 
\begin{equation}
	\label{eq:mismatches}
	\begin{aligned}
		\Delta P_k &= P_k^{\text{spec}} - P_k^{\text{calc}} \\
		\Delta Q_k &= Q_k^{\text{spec}} - Q_k^{\text{calc}}
	\end{aligned}
\end{equation}
werden als Zustandsgrößen verwendet \todo{literatur.}.
Die Berechnung der Leistungsgrößen $P_k^{\text{calc}}$ und $Q_k^{\text{calc}}$ geschieht über die Auswertung des Produkts zwischen Spannung und Strom \eqref{eq:power_current}.
Ausformuliert folgt
%als Differenz zwischen vorgegebenen und aus den aktuellen Zustandsgrößen berechneten Leistungen.
%Die berechneten Leistungen 
\begin{equation}
	\label{eq:pq_calc}
	\begin{aligned}
		P_k^{\text{calc}} &= |V_k| \sum_{m \in \Omega_B} |V_m| \left( G_{km}\cos\theta_{km} + B_{km}\sin\theta_{km} \right) \\
		Q_k^{\text{calc}} &= |V_k| \sum_{m \in \Omega_B} |V_m| \left( G_{km}\sin\theta_{km} - B_{km}\cos\theta_{km} \right)
	\end{aligned}
\end{equation}
unter Verwendung von \eqref{eq:y_split} und der Winkeldifferenz $\theta_{km} \equiv \theta_k - \theta_m$.
Die Summation erfolgt über alle Knoten $m$, da über die Nebendiagonalelemente $Y_{km}$ alle mit $k$ gekoppelten Knoten Beiträge zur Leistungseinspeisung am Knoten $k$ liefern. \todo{literatur.} \\


%Das nichtlineare Gleichungssystem $\mathbf{f}(\mathbf{x}) = \mathbf{0}$ mit Zustandsvektor $\mathbf{x} = [\boldsymbol{\theta}^\top, |\mathbf{V}|^\top]^\top$ wird durch Taylor-Entwicklung erster Ordnung um den aktuellen Iterationspunkt linearisiert:

%Das NR-Verfahren linearisiert den Mismatch-Vektor 
Im NR-Verfahren wird die Leistungsdifferenz (engl. \emph{power mismatch}) durch den Vektor 

\begin{equation}
	\vect{g}(\vect{x}) = 
	\begin{bmatrix}
		\Delta \vect{P} \\
		\Delta \vect{Q}
	\end{bmatrix}
	= \vect{0}
	\label{eq:mismatch_polar}
\end{equation}

um den aktuellen Arbeitspunkt
 \begin{equation}
   \vect{x}^{(k)}= \begin{bmatrix}
   	\boldsymbol{\theta} \\
   	 |\mathbf{V}|
   \end{bmatrix}
 	\label{eq:arbeitspunkt}
 \end{equation}
durch eine Taylor-Reihen-Entwicklung erster Ordnung linearisiert.
Eine kleine Deviation um den Arbeitspunkt kann folglich mit
\begin{equation}
	\vect{g}(\vect{x}^{(k)} + \Delta \vect{x}) \approx \vect{g}(\vect{x}^{(k)}) + \vect{J}(\vect{x}^{(k)}) \cdot \Delta \vect{x} = \vect{0}
	\label{eq:nr_taylor}
\end{equation}
dargestellt werden. 
Ergänzend ist anzumerken, dass  $\vect{J}$ die Jacobi-Matrix ist, die folgend näher betrachtet wird.
Genau hier kann nach der Korrektur
\begin{equation}
	\Delta \vect{x}^{(k)} = -\vect{J}(\vect{x}^{(k)})^{-1} \cdot \vect{g}(\vect{x}^{(k)})
	\label{eq:newton_richtung}
\end{equation}
aufgelöst werden.
Daraus kann letztlich die Iterationsvorschrift 
\begin{equation}
	\vect{x}^{(k+1)} = \vect{x}^{(k)} + \Delta \vect{x}^{(k)}
	\label{eq:nr_update}
\end{equation}
formuliert werden, die den Zustandsvektor bei jeder Iteration aktualisiert \todo{literatur.}.\\ 

%
%Ausgeschrieben bedeutet dies:
%\begin{align}
%	\theta_k^{(k+1)} &= \theta_k^{(k)} + \Delta \theta_k^{(k)} \quad &&\text{für } k \in \Omega_{PQ} \cup \Omega_{PV} \label{eq:update_theta}\\
%	|V_k|^{(k+1)} &= |V_k|^{(k)} + \Delta |V_k|^{(k)} \quad &&\text{für } k \in \Omega_{PQ} \label{eq:update_V}
%\end{align}
%
%Der Algorithmus wird initialisiert mit dem \textit{Flat Start}: $|V_k|^{(0)} = 1{,}0 \pu$ und $\theta_k^{(0)} = 0$ für alle Knoten (außer Slack). Die Iteration wird beendet, wenn $\max(|\Delta P_k|, |\Delta Q_k|) < \varepsilon$ mit typischen Toleranzen $\varepsilon = 10^{-6}$ bis $10^{-8}$.
%Die Jacobi-Matrix $\vect{J}$ enthält die partiellen Ableitungen der Mismatch-Funktionen nach den Unbekannten und wird in der polaren Formulierung in vier Blöcke unterteilt \cite{Grainger1994}:

Die Jacobi-Matrix umfasst die partiellen Differenziationen der Leistungsdifferenzen nach den einzelnen Zustandsgrößen \todo{literatur.}.
Durch die eingeführte Struktur von zwei Leistungs- und Spannungsgrößen können die Jacobi-Einträge in vier Blöcke 
\begin{equation}
	\mathbf{H} = \frac{\partial \mathbf{P}}{\partial \boldsymbol{\theta}}, \quad
	\mathbf{N} = \frac{\partial \mathbf{P}}{\partial |\mathbf{V}|}, \quad
	\mathbf{M} = \frac{\partial \mathbf{Q}}{\partial \boldsymbol{\theta}}, \quad
	\mathbf{L} = \frac{\partial \mathbf{Q}}{\partial |\mathbf{V}|}
\end{equation}
gegliedert werden \todo{literatur.}. 
Zusammengesetzt und eingesetzt in Gleichung \eqref{eq:nr_taylor} ist der Zusammenhang der Ableitungen mit den gesuchten und bekannten Größen durch 
%Die Blöcke der Jacobi-Matrix $\mathbf{J} \in \mathbb{R}^{2b \times 2b}$ sind definiert als

%\begin{equation}
%	\begin{bmatrix}
%		\Delta \vect{P} \\
%		\Delta \vect{Q}
%	\end{bmatrix}
%	=
%	\underbrace{
%		\begin{bmatrix}
%			\vect{H} & \vect{N} \\
%			\vect{M} & \vect{L}
%		\end{bmatrix}
%	}_{\vect{J}}
%	\begin{bmatrix}
%		\Delta \boldsymbol{\theta} \\
%		\Delta |\vect{V}|
%	\end{bmatrix}
%	\label{eq:jacobi_block}
%\end{equation}


\begin{equation}
	\label{eq:nr_linearization}
	\underbrace{
	\begin{bmatrix} \Delta \mathbf{P} \\ \Delta \mathbf{Q} \end{bmatrix}
	}_{\mathbf{g}(\mathbf{x})}
	=
	\underbrace{
		\begin{bmatrix}
			\mathbf{H} & \mathbf{N} \\
			\mathbf{M} & \mathbf{L}
		\end{bmatrix}
	}_{\mathbf{J}}
	\begin{bmatrix} \Delta \boldsymbol{\theta} \\ \Delta |\mathbf{V}|\end{bmatrix}
\end{equation}

hergestellt, wobei die Jacobi dementsprechend die Dimension $\mathbf{J} \in \mathbb{R}^{2b \times 2b}$ hat.
Damit kann die Iteration durch Aktualisierung der Zustandsgrößen, Jacobi, und Leistungsgrößen durchgeführt werden. 
%Die Iterationsvorschrift lautet
%\begin{equation}
%	\label{eq:nr_update}
%	\mathbf{x}^{(n+1)} = \mathbf{x}^{(n)} + \mathbf{J}^{-1}(\mathbf{x}^{(n)}) \cdot \Delta \mathbf{f}(\mathbf{x}^{(n)})
%\end{equation}
Diese wird durchgeführt, solange eine gewählte Konvergenzgrenze $\varepsilon$ durch
\begin{equation}
	\label{eq:convergence_limit}
	\max(\|\Delta \mathbf{P}\|, \|\Delta \mathbf{Q}\|) < \varepsilon
\end{equation}
nicht erreicht wird.\\

%Der wesentliche Vorteil des NR liegt in seiner \textbf{quadratischen Konvergenz} in der Nähe der Lösung. In praktischen Anwendungen konvergiert das Verfahren typischerweise in 3--5 Iterationen~\cite{Salazar2024}. Die Kosten pro Iteration setzen sich zusammen aus dem Aufbau der Jacobi-Matrix, der LU-Faktorisierung des Sparse-Systems \cite{Davis2006}. Der Hauptnachteil ist, dass \emph{beide} Schritte in jeder Iteration wiederholt werden müssen, da die Jacobi-Matrix vom aktuellen Zustand $\mathbf{x}^{(n)}$ abhängt.


\subsection{Behandlung von PV-Knoten}
\label{subsec:nr_pv}

Die bisherige Einführung des Verfahrens stützt sich auf Knoten wo Wirk- und Blindleistung bekannt sind, also PQ-Knoten.
PV-Knoten liefern dahingehend nur die Wirkleistungen $|P_k^{\text{spec}}|$ aus dem Leistungsvektor und die Spannungsbeträge 
$|V_k^{\text{spec}}|$ aus dem Zustandsvektor, während die Blindleistung $Q_k$ und der Winkel $\theta_k$ gesucht bleiben.
Daraus folgt eine Modifikation des NR-Verfahrens zur korrekten Behandlung dieser Knoten.\\

%An PV-Knoten $k \in \Omega_{PV}$ ist der Spannungsbetrag $|V_k| = |V_k^{\text{spec}}|$ fest vorgegeben, während die Blindleistung $Q_k$ eine unbekannte Größe darstellt. Dies erfordert eine Modifikation des Standard-NR.

Die verbreitete Implementierung, wie sie auch in der Lastflussberechnungssoftware \texttt{pandapower} vorzufinden ist, erfordert eine Systemreduktion.
Die Blindleistungsgleichungen für PV-Knoten werden entfernt und der Spannungsbetrag wird als gegebener Wert festgehalten.
Durch das Weglassen der Zeilen und Spalten wird die Dimension auf des Systems auf $(2n_{PQ} + n_{PV})$ verkleinert.
Die Lastflussgleichung mit ihren Jacobi-Blöcken ist nach Reduktion durch
%reduziert die Systemgröße, indem für jeden PV-Knoten die Q-Gleichung aus dem Gleichungssystem entfernt wird und die zugehörige Spannungsvariable $|V_k|$ als konstant fixiert wird. Das reduzierte System hat die Dimension $(2n_{PQ} + n_{PV})$ statt $2b$:
\begin{equation}
	\label{eq:nr_reduziert}
	\begin{bmatrix} \Delta \mathbf{P} \\ \Delta \mathbf{Q}_{\mathrm{PQ}} \end{bmatrix}
	=
	\begin{bmatrix}
		\mathbf{H} & \mathbf{N}_{\mathrm{PQ}} \\
		\mathbf{M}_{\mathrm{PQ}} & \mathbf{L}_{\mathrm{PQ,PQ}}
	\end{bmatrix}
	\begin{bmatrix} \Delta \boldsymbol{\theta} \\ \Delta |\mathbf{V}_{\mathrm{PQ}}|\end{bmatrix}
\end{equation}

\begin{equation}
	\mathbf{H} = \frac{\partial \mathbf{P}}{\partial \boldsymbol{\theta}}, \quad
	\mathbf{N}_{PQ} = \frac{\partial \mathbf{P}}{\partial |\mathbf{V}_{\mathrm{PQ}}|}, \quad
	\mathbf{M}_\mathrm{PQ} = \frac{\partial \mathbf{Q}_{\mathrm{PQ}}}{\partial \boldsymbol{\theta}}, \quad
	\mathbf{L}_{\mathrm{PQ,PQ}} = \frac{\partial \mathbf{Q}_{\mathrm{PQ}}}{\partial |\mathbf{V}_{\mathrm{PQ}}|}
\end{equation}

gegeben.
Nach vollständiger Konvergenz werden die Blindleistungen der PV-Knoten aus \eqref{eq:pq_calc} nachträglich berechnet.\\


Dabei ist der Nachteil des Verfahrens die Neuberechnung und Aufbau der Jacobi-Matrix in jedem Iterationsschritt, da die Ableitungen direkt von dem aktuellen Zustand abhängen~\cite{Davis2006}. 
Den Vorteil hingegen bildet die quadratische Konvergenzeigenschaft der Verfahrens. Dadurch findet das Verfahren im Normalfall zwischen 3--5 Iterationen eine Lösung~\cite{Salazar2024}.

%Ein alternativer Ansatz formuliert das NR-Verfahren nicht in Polar-, sondern in \emph{kartesischen Koordinaten} auf Basis von \emph{Stromeinspeisungs-Ungleichungen} (\emph{current injection mismatches}) $\Delta I_{r,k}$ und $\Delta I_{m,k}$ (Real- und Imaginärteil). Für PQ-Knoten sind die Nichtdiagonalblöcke der resultierenden $(2n \times 2n)$-Jacobi-Matrix identisch mit den entsprechenden $(2 \times 2)$-Blöcken der Knotenadmittanzmatrix und bleiben über alle Iterationen konstant; nur die Diagonalblöcke müssen aktualisiert werden.














%strominjektion



%In einer alternativen Formulierung können PV-Knoten anderweitig im NR-Verfahren eingefügt werden. 
%Nach Costa et al. \cite{daCosta1999} kann das Verfahren in kartesischen Koordinaten basierend auf realen und imaginären Stromeinspeisungen $\Delta I_{r,k}$ und $\Delta I_{m,k}$ formuliert werden.
%Anzumerken ist, dass die Nebendiagonalblöcke der PQ-Knoten der Jacobimatrix mit den Knotenadmittanzmatrizen übereinstimmen und deswegen über alle Iterationen konstant bleiben.
%
%\todo{hier weitermachen}
%An einem PV-Knoten $k$ wird die Blindleistung $\Delta Q_k$ als zusätzliche abhängige Variable eingeführt und die zusätzliche Gleichung $\Delta |V_k| = 0$ hinzugefügt, sodass zunächst ein erweitertes System der Dimension $(2n + n_{PV})$ entsteht:
%\begin{equation}
%	\label{eq:nr_erweitert}
%	\begin{bmatrix}
%		\Delta I_{m,k} \\
%		\Delta I_{r,k} \\
%		\Delta |V_k|
%	\end{bmatrix}
%	=
%	\begin{bmatrix}
%		B_{kk} & G_{kk} & \partial I_{m,k}/\partial Q_k \\
%		G_{kk} & -B_{kk} & \partial I_{r,k}/\partial Q_k \\
%		\partial |V_k|/\partial V_{r,k} & \partial |V_k|/\partial V_{m,k} & 0
%	\end{bmatrix}
%	\begin{bmatrix}
%		\Delta V_{r,k} \\
%		\Delta V_{m,k} \\
%		\Delta Q_k
%	\end{bmatrix}
%\end{equation}
%Die letzte Zeile von \eqref{eq:nr_erweitert} liefert die lineare Beziehung $\Delta V_{m,k} = -(V_{r,k}/V_{m,k})\,\Delta V_{r,k}$, die in die ersten beiden Zeilen eingesetzt wird. Dadurch wird eine der Spannungsvariablen eliminiert und durch $\Delta Q_k$ ersetzt, sodass das System auf die ursprüngliche Dimension $(2n \times 2n)$ \emph{reduziert} wird. Der wesentliche Vorteil liegt darin, dass die Jacobi-Matrix unabhängig von der Anzahl der PV-Knoten die Struktur der Knotenadmittanzmatrix beibehält: Die Nichtdiagonalblöcke bleiben konstant, und nur wenige Diagonalblöcke müssen pro Iteration aktualisiert werden.\\
%Empirische Ergebnisse an praktischen Netzen zeigen, dass dieser Ansatz die Gesamtrechenzeit gegenüber der klassischen NR-Formulierung in Polarkoordinaten um ca.\ 20\,\% reduziert, bei identischem Konvergenzverhalten~\cite{daCosta1999}. In dieser Arbeit wird dennoch die polare Standardformulierung als NR-Referenz verwendet, da diese der Implementierung in \texttt{pandapower} entspricht, welches als Referenz dient.

%\paragraph{Behandlung von Q-Grenzen.}
%In der Praxis sind die an einem PV-Knoten einspeisbaren Blindleistungen durch $Q_k^{\min} \leq Q_k \leq Q_k^{\max}$ begrenzt (vgl.\ \eqref{eq:qlimits}). Nach jeder NR-Iteration wird die berechnete Blindleistung mit den Grenzen verglichen. Wird eine Grenze verletzt, so wird der Knoten in einen PQ-Knoten mit $Q_k = Q_k^{\min}$ bzw.\ $Q_k = Q_k^{\max}$ umgewandelt und die Spannungsbedingung fallengelassen. Umgekehrt wird ein zuvor umgeschalteter Knoten wieder in einen PV-Knoten überführt, sobald die Spannungsbedingung wieder erfüllbar wird. Dieser \emph{Q-Limit-Enforcement}-Mechanismus wird auch in den vorgeschlagenen Tensor-Methoden (Kap.~\ref{sec:pv_methoden}) übernommen.

\subsection{Bewertung für multidimensionale Lastflussanalysen}
\label{subsec:nr_massen}

Genau dieser Vor- und Nachteil soll im Bezug auf multidimensionale Berechnungen näher betrachtet werden.
Bei mehreren Lastflussszenarien $\tau$, wie sie in der Zeitreihenberechnung, Monte-Carlo-Simulationen und probabilistischen Lastflüssen vorkommen, folgt ein struktureller Nachteil: Die Jacobi-Matrix $\mathbf{J}(\mathbf{x}^{(n)}, \mathbf{s})$ muss aufgrund Ihrer Abhängigkeit von der Leistung und Zustand nicht nur in jeder Iteration bestimmt werden, sondern auch in jedem Szenario.\\
%Die 
%
%
%Für einzelne Lastflüsse ist das NR-Verfahren aufgrund seiner quadratischen Konvergenz und der ausgereiften Sparse-Implementierungen sehr effizient. Für die in dieser Arbeit betrachteten (\emph{multidimensionalen}) Anwendungen, probabilistische Lastflüsse und Sensitivitätsstudien mit $\tau \gg 1$ Szenarien, ergibt sich jedoch ein struktureller Nachteil:
%\begin{itemize}
%	\item Die Jacobi-Matrix $\mathbf{J}(\mathbf{x}^{(n)}, \mathbf{s})$ hängt sowohl vom Iterationszustand als auch von den szenarienspezifischen Leistungsvorgaben ab. Sie muss daher in jeder Iteration \emph{jedes} Szenarios neu aufgebaut und faktorisiert werden.
%	\item Der Gesamtaufwand skaliert linear mit $\tau$: $t_{\text{gesamt}} \approx \tau \cdot t_{\text{NR}}$.
%	\item Eine Parallelisierung über Szenarien ist zwar prinzipiell möglich, aber ineffizient, da jede Instanz eine eigene Sparse-Faktorisierung durchführen muss. Standard-BLAS-Bibliotheken (\texttt{OpenBLAS}, \texttt{MKL}) können nicht effektiv genutzt werden.
%\end{itemize}

%Empirisch wird anhand einer Jahressimulation im Minutenraster ($\tau = 525\,600$ Lastflüsse) für ein 100-Bus-Netz folgende Rechenzeiten belegt~\cite{Salazar2024}:

In Tabelle \ref{tab:nr_vergleich} ist exemplarisch ein Ergebnis für ein 100-Knoten-Netz einer Jahressimulation abgebildet~\cite{Salazar2024}, welches diesen Nachteil belegt. 
\begin{table}[H]
	\centering
	\caption{Rechenzeit in Minuten für verschiedene Anzahl an Lastflüsse (PF) eines 100-Knoten-Netzes~\cite{Salazar2024}}
	\label{tab:nr_vergleich}
\begin{tabular}{llrrrr}
	\toprule
	Algorithmus & Netzgröße &
	\SI{8760}{PF} &
	\SI{17520}{PF} &
	\SI{35040}{PF} &
	\SI{525600}{PF} \\
	\midrule
	NR
	& 100
	& 1.03
	& 2.06
	& 4.12
	& 61.85 \\
	
	Tensor
	& 100
	& 0.01
	& 0.01
	& 0.02
	& 0.37 \\

	\bottomrule
\end{tabular}
\end{table}

Dieses Ergebnis begründet den nahezu linearen Anstieg der Rechenzeit bei ansteigender Szenarienanzahl $t_{\text{gesamt}} \approx \tau \cdot t_{\text{NR}}$.
Dabei konnte bereits für den reinen PQ-Knoten-Fall gezeigt werden dass eine tensorisierte Fixpunktberechnung um den Faktor 164 schneller sein kann als das NR-Verfahren~\cite{Salazar2024}. Dies liegt, wie folgend erklärt an der Konstanz der Systemmatrix in derm Fixpunktverfahren gegenüber der sich verändernden Jacobimatrix.
%Der Faktor 164 zwischen NR (Sparse) und Tensor (Dense) verdeutlicht den Kern der Motivation dieser Arbeit. Bei fester Netztopologie und konstantem Leistungsmodell bleibt die zentrale Systemmatrix der Fixpunkt-Iteration ($\mathbf{Z}_B$) über alle Szenarien und Iterationen \emph{konstant}. Die dadurch mögliche Vorab-Berechnung und -Faktorisierung sowie die anschließende Batched-Matrix-Multiplikation eliminieren den Hauptaufwand pro Szenario. Diese strukturelle Eigenschaft, die dem NR fundamental fehlt, ist der Ausgangspunkt für die im nächsten Abschnitt vorgestellte Fixpunkt-Iteration.

\bigskip
\noindent
%Zusammenfassend wurde das NR-Verfahren mit seiner Standardformulierung in Polarkoordinaten und den beiden gängigen Ansätzen zur Behandlung von PV-Knoten dargestellt. Trotz seiner exzellenten Eigenschaften pro Einzellastfluss ist es aufgrund der szenarienabhängigen Jacobi-Matrix nicht für massenparallele Berechnungen geeignet. Dies motiviert die im folgenden Abschnitt behandelte Fixpunkt-Iteration mit konstanter Systemmatrix.


\section{Fixpunkt-Iteration und Banachscher Fixpunktsatz}
\label{sec:fpi}

%Dieses Kapitel führt die Fixpunkt-Iteration (FPI) als alternatives Lösungsverfahren für das Lastflussproblem ein und stellt damit das mathematische Fundament des in Kapitel~\ref{sec:tpf} entwickelten Tensor Power Flow bereit. Im Gegensatz zum Newton-Raphson-Verfahren verzichtet die FPI auf eine Linearisierung und arbeitet stattdessen mit einer expliziten Iterationsvorschrift, deren zentrale Systemmatrix über alle Iterationen \emph{konstant} bleibt. Diese Eigenschaft ist der Ausgangspunkt für die spätere Massenparallelisierung.


Die Fixpunkt-Iteration (FPI) wird in diesem Abschnitt eingeführt. 
Das Verfahren dient als tensorisierbare Alternative zu dem Newton-Raphson-Verfahren.
Dabei wird gezeigt, dass für Netze ohne PV-Knoten die Systemmatrix über alle Iterationen konstant bleibt, nicht wie die Jacobimatrix in dem NR-Verfahren.
Dieses Kapitel dient weiterhin als Basis für die Parallelisierung, welche im \cref{ch:tpf} des \emph{Tensor Power Flow} durchgeführt wird.


%Die Darstellung folgt der Herleitung von Giraldo et al.~\cite{Giraldo2022} über die Laurent-Reihen-Entwicklung und der äquivalenten Formulierung in Salazar Duque et al.~\cite{salazar2024}. Abschnitt~\ref{subsec:fpi_herleitung} leitet die Iterationsvorschrift her, Abschnitt~\ref{subsec:banach} beweist die Konvergenz über den Banachschen Fixpunktsatz, Abschnitt~\ref{subsec:fpi_interpretation} liefert die physikalische Interpretation und Abschnitt~\ref{subsec:fpi_eigenschaften} fasst die für die Tensorformulierung relevanten Eigenschaften zusammen.



\subsection{Laurent-Reihen-Approximation}
\label{subsec:laurent}


In der Leistungsgleichung \eqref{eq:power_current} ist ersichtlich dass durch die elementweise Division der Spannung $V^{-1}$ eine Nichtlinearität vorliegt.
Diese pflanzt sich in die Darstellung des leistungsabhängigen Knotenstromes fort~\cite{Giraldo2022}.
Um diese zu linearisieren, wird eine Laurent-Reihenapproximation durchgeführt.
Dabei wird die spannungsabhängige Funktion $f(V) = V^{-1}$ an einem Arbeitspunkt $\mathrm{V}^{(0)}$ als eine Laurent-Reihe erster Ordnung

\begin{equation}
	f(V_i) \approx \frac{2}{\mathrm{V}^0} - \frac{V_i}{(\mathrm{V}^0)^2}
	\label{eq:laurent_approx}
\end{equation}

dargestellt werden, wobei
 
\begin{equation}
	V_i = \mathrm{V}^0 - \Delta V \quad \text{mit} \quad \|\Delta V\| / \|\mathrm{V}^0\| < 1
	\label{eq:annahme}
\end{equation}
 angenommen wird.
%Die Nichtlinearität der Lastflussgleichungen resultiert aus dem Term $(V_k)^{-1}$, der im spezifizierten Strom \eqref{eq:power_current} auftritt. Giraldo et al.\ \cite{Giraldo2022} schlagen eine Linearisierung mittels Laurent-Reihenentwicklung vor.
%	Die Nichtlinearität $f(V_i) = (V_i)^{-1}$ an einem Knoten $i \in \Omega_d$ wird um einen Arbeitspunkt $\mathrm{V}^0 \in \C$ mittels Laurent-Reihe approximiert als:
%	\begin{equation}
%		f(V_i) \approx \frac{2}{\mathrm{V}^0} - \frac{V_i}{(\mathrm{V}^0)^2}
%		\label{eq:laurent_approx}
%	\end{equation}
%	unter der Annahme $V_i = \mathrm{V}^0 - \Delta V$ mit $\|\Delta V\| / \|\mathrm{V}^0\| < 1$.
%\begin{proof}
%	Die Funktion $f(\Delta V) = 1/(\mathrm{V}^0 - \Delta V)$ kann für $\|\Delta V\| < \|\mathrm{V}^0\|$ als konvergente Potenzreihe dargestellt werden \cite{Giraldo2022}:
%	\begin{equation}
	%		f(\Delta V) = \sum_{n=1}^{\infty} \frac{\Delta V^{n-1}}{(\mathrm{V}^0)^n} = \frac{1}{\mathrm{V}^0} + \frac{\Delta V}{(\mathrm{V}^0)^2} + \ldots
	%		\label{eq:laurent_reihe}
	%	\end{equation}
%	Durch Abbruch nach dem zweiten Glied ($n \leq 2$) und Rücksubstitution $\Delta V = \mathrm{V}^0 - V_i$ ergibt sich die Approximation \eqref{eq:laurent_approx}.
%\end{proof}
Diese Annahme wird dadurch begründet, dass die Spannungen in Verteilernetzen innerhalb eines engen Bandes um den Nennwert liegen. 
Das Band erstreckt sich hierbei typischerweise von 0,9 bis $1,1\,\mathrm{p.u.}$~\cite{Giraldo2022}.
%Im iterativen Algorithmus wird der Arbeitspunkt $\mathrm{V}^0$ nach jeder Iteration auf den aktuellen Spannungswert aktualisiert, sodass die Approximation mit jeder Iteration genauer wird.

\subsection{Herleitung der Fixpunkt-Iteration}
\label{subsec:fpi_herleitung}

Das ZIP-Modell \eqref{eq:zip} kann durch die Knotenspannung $V_\mathrm{d}$ dividiert werden um die Knotenströme 

\begin{equation}
	\label{eq:current_zip}
	%		-\operatorname{diag}(\mathbf{v}_d^{*})^{-1}\mathbf{s}_d^{*} &= \mathbf{Y}_{ds}\mathbf{v}_s + \mathbf{Y}_{dd}\mathbf{v}_d\\
	%	&\Updownarrow\\
	\mathbf{i}_d = \left(\boldsymbol{\alpha}_P \odot \operatorname{diag}(\mathbf{v}_d^{*})^{-1} + \operatorname{diag}(\boldsymbol{\alpha}_I) + \boldsymbol{\alpha}_Z \odot \operatorname{diag}(\mathbf{v}_d)\right)\mathbf{s}_n^{*}
\end{equation}

als Funktion der ZIP-aufgeteilten Spannung und Nennleistung darzustellen.
An dieser Stelle ist die Nichtlinearität $\operatorname{diag}(\mathbf{v}_d^{*})^{-1}$ sichtbar. 
Für diesen Ausdruck wird die bereits eingeführte Laurent-Reihe verwendet. 
Der Ausdruck für den Knotenstrom wird weiterhin in die Knotenbilanz \eqref{eq:current_balance} eingesetzt und führt zu



%Ausgangspunkt ist die Grundgleichung des Lastflusses \eqref{eq:pf_fundamental}. Unter Einbeziehung des ZIP-Lastmodells \eqref{eq:zip} lässt sich die Nichtlinearität in eine explizite Iterationsvorschrift überführen. Nach Giraldo et al.~\cite{Giraldo2022} ergibt sich die Stromgleichung an den Lastknoten zu

%Einsetzen in die Knotenstrombilanz \eqref{eq:current_balance} und Umordnen führt auf das lineare System
\begin{equation}
	\vect{A}\vect{V}_d^* - \vect{B}\vect{V}_d = \vect{C} + \vect{D}
	\label{eq:lineares_system}
\end{equation}
mit den Matrizen
\begin{align}
	\vect{A} &= \diag\left(\boldsymbol{\alpha_P} \had \mathrm{V}^{0*-2} \had \mathbf{S}_n^*\right) \label{eq:giraldo_A}\\
	\vect{B} &= \diag(\boldsymbol{\alpha_Z} \had \mathbf{S}_n^*) + \mathbf{\Ydd} \label{eq:giraldo_B}\\
	\vect{C} &= \mathbf{\Yds} \mathbf{V}_s + \boldsymbol{\alpha_I} \had \mathbf{S}_n^* \label{eq:giraldo_C}\\
	\vect{D} &= 2\boldsymbol{\alpha_P} \had \mathrm{V}^{0*-1} \had \mathbf{S}_n^* \label{eq:giraldo_D}.
\end{align}
Auflösen von \eqref{eq:lineares_system} nach der Spannung $\mathbf{V}_d$ liefert schließlich die Iterationsvorschrift des Fixpunkts
\begin{equation}
	\label{eq:fpi_iteration}
	\boxed{\;	\vect{V}_d^{(k+1)} = \vect{B}^{-1}\left(\vect{A}^{(k)} \vect{V}_d^{*(k)} - \vect{C} - \vect{D}^{(k)}\right) =: T\!\left(\mathbf{v}_d^{(n)}\right)\;}
\end{equation}
nach Giraldo et al.~\cite{Giraldo2022}.
Ähnlich zu dem NR-Verfahren wird diese Rechnung bis zum erreichen der Konvergenzgrenze $\|\mathbf{v}_d^{(n+1)} - \mathbf{v}_d^{(n)}\| < \varepsilon$ wiederholt.
%Die Iteration wird typischerweise mit einem Flat-Start $\mathbf{v}_d^{(0)} = \mathbf{1}$ initialisiert und abgebrochen, sobald $\|\mathbf{v}_d^{(n+1)} - \mathbf{v}_d^{(n)}\| < \varepsilon$.
Diese Darstellung ist für jegliche ZIP-Lastkonfiguration gültig. 

%Um die Abhängigkeit der Leistung in dem nachfolgenden Kapitel besser Aufzeigen zu können
%%Für das in dieser Arbeit als Standardfall angenommene konstante Leistungsmodell ($\boldsymbol{\alpha}_P = \mathbf{1}$, $\boldsymbol{\alpha}_I = \boldsymbol{\alpha}_Z = \mathbf{0}$) vereinfachen sich \eqref{eq:fpi_ABc} zu
%\begin{equation}
%	\label{eq:fpi_konstant}
%	\mathbf{A} = \operatorname{diag}(\mathbf{s}_d^{*}), \quad \mathbf{B} = \mathbf{Y}_{dd}, \quad \mathbf{c} = \mathbf{Y}_{ds}\mathbf{v}_s
%\end{equation}
%Mit der Bus-Impedanzmatrix $\mathbf{Z}_B = \mathbf{Y}_{dd}^{-1}$ aus \eqref{eq:zbus} lässt sich die Iteration als
%\begin{equation}
%	\label{eq:fpi_zbus}
%	\mathbf{v}_d^{(n+1)} = \mathbf{Z}_B\left(\mathbf{s}_d^{*} \oslash \mathbf{v}_d^{*\,(n)}\right) - \mathbf{Z}_B\mathbf{Y}_{ds}\mathbf{v}_s
%\end{equation}
%schreiben. Diese kompakte Form ist die Grundlage der Tensorformulierung, da $\mathbf{Z}_B$ und $\mathbf{Z}_B\mathbf{Y}_{ds}\mathbf{v}_s$ szenarienunabhängig konstant sind und einmalig vorab berechnet werden können.
%
%
%\begin{algorithm}[H]
%	\caption{Successive Approximation Method (SAM) nach \cite{Giraldo2022}}
%	\label{alg:sam}
%	\begin{algorithmic}[1]
%		\REQUIRE $\vect{S}_n$, $\Ydd$, $\Yds$, $\vect{V}_s$, Toleranz $\varepsilon$, max.\ Iterationen $K$
%		\ENSURE $\vect{V}_d$
%		
%		\STATE Initialisiere $k = 0$, $\text{tol} = \infty$
%		\STATE Initialisiere $\mathrm{V}^0_{i,\phi} = [1, a^2, a]$ für alle $i \in \Omega_d$ \COMMENT{Flat Start, $a = e^{\jj 2\pi/3}$}
%		\STATE Berechne $\vect{B}^{-1}$ und $\vect{C}$ gemäß \eqref{eq:giraldo_B}--\eqref{eq:giraldo_C}
%		
%		\WHILE{$\text{tol} \geq \varepsilon$ \textbf{und} $k < K$}
%		\STATE Berechne $\vect{A}^{(k)}$ und $\vect{D}^{(k)}$ gemäß \eqref{eq:giraldo_A}--\eqref{eq:giraldo_D}
%		\STATE $\vect{V}_d^{(k+1)} = \vect{B}^{-1}\left(\vect{A}^{(k)} \vect{V}_d^{*(k)} - \vect{C} - \vect{D}^{(k)}\right)$
%		\STATE Berechne Leistungsmismatch: $\vect{S}_d^{*(k+1)} = \diag(\vect{V}_d^{*(k+1)})(\Yds\vect{V}_s + \Ydd\vect{V}_d^{(k+1)})$
%		\STATE $\text{tol} \leftarrow \max\|\vect{S}_d^* - \vect{S}_d^{*(k+1)}\|$
%		\STATE $\mathrm{V}^0_d \leftarrow \vect{V}_d^{(k+1)}$ \COMMENT{Arbeitspunkt aktualisieren}
%		\STATE $k \leftarrow k + 1$
%		\ENDWHILE
%		
%		\STATE Berechne $\vect{S}_s^* = \diag(\vect{V}_s^*)(\vect{Y}_{ss}\vect{V}_s + \vect{Y}_{sd}\vect{V}_d^{(k)})$
%		\RETURN $\vect{V}_d^{(k)}$, $\vect{S}_s$
%	\end{algorithmic}
%\end{algorithm}


\subsection{Konvergenz über den Banachschen Fixpunktsatz}
\label{subsec:banach}

%Die Konvergenz der Iteration \eqref{eq:fpi_iteration} lässt sich über den Banachschen Fixpunktsatz beweisen. Dieser besagt, dass eine kontrahierende Abbildung $T: X \to X$ auf einem vollständigen metrischen Raum $(X,d)$ genau einen Fixpunkt $\mathbf{u} = T(\mathbf{u})$ besitzt, zu dem die Folge $\mathbf{v}^{(n+1)} = T(\mathbf{v}^{(n)})$ für jeden Startwert $\mathbf{v}^{(0)} \in X$ konvergiert. Die Abbildung heißt \emph{kontrahierend} mit Kontraktionskonstante $\eta$, wenn
%\begin{equation}
%	\label{eq:contraction}
%	d\!\left(T(\mathbf{v}),\, T(\mathbf{w})\right) \leq \eta\, d(\mathbf{v}, \mathbf{w}) \quad \forall \mathbf{v},\mathbf{w} \in X, \quad 0 \leq \eta < 1
%\end{equation}

Die Konvergenz des Verfahrens ist über den Banachschen Fixpunktsatz charakterisiert \cite{Banach.1922}. Dieser besagt eine kontrahierende Abbildung $T: X \to X$ einen Fixpunkt $u$ besitzt, zudem jeder Startpunkt konvergiert.
Übertragen auf das Lastflussproblem ist $\mathbf{u}$ die exakte Lösung des Lastflusses. Giraldo et al.~\cite{Giraldo2022} wendet der Fixpunktsatz auf den FPI an leitet die Ungleichung
\begin{equation}
	\label{eq:contraction_tpf}
	\left\|\mathbf{v}_d^{(n+1)} - \mathbf{u}_d\right\|
	\leq \frac{\left\|\mathbf{B}^{-1}\operatorname{diag}(\boldsymbol{\alpha}_P \odot \mathbf{s}_n^{*})\right\|}{v_{\min}^2}
	\left\|\mathbf{v}_d^{(n)} - \mathbf{u}_d\right\|
\end{equation}

her, wobei $v_{\min}$  der minimale Spannungsbetrag im Netz ist. Nach Banach ist der Vorfaktor 

\begin{equation}
	\label{eq:eta}
	\boxed{\;\eta = \max_{k \in \Omega_d}\left\{\frac{\left\|\mathbf{B}^{-1}\operatorname{diag}(\boldsymbol{\alpha}_P \odot \mathbf{s}_n^{*})\right\|}{v_{\min}^2}\right\}\;}
\end{equation}

die Kontraktionskonstante des Systems.
Die FPI konvergiert genau dann für einen beliebigen Startwert zur eindeutigen Lösung, wenn $\eta < 1$ gilt. 
Ferner ist zu beachten, dass dies nur eine hinreichende Bedingung ist und für  $\eta > 1$ die Iteration trotzdem konvergieren kann, diese jedoch nicht mehr garantiert ist.\\

Zur besseren Interpretation der Konstante wird der konstante Lastfall $\boldsymbol{\alpha}_\mathrm{P} = \mathbf{1}$ verwendet. 
Dabei kann unter der Verwendung von $\mathbf{z}_{L,k} = v_{\min}^2 / |s_{n,k}^{*}|$ als Lastimpedanz an Knoten $k$ die Kontraktionskonstante auf 
%Als Interpretation für den Fall konstanter Leistung ($\boldsymbol{\alpha}_P = \mathbf{1}$) und unter Verwendung der Identität $\mathbf{s}_n^{*} \oslash v_{\min}^2 = \mathbf{1} \oslash \mathbf{z}_{L}$, wobei $\mathbf{z}_{L,k} = v_{\min}^2 / |s_{n,k}^{*}|$ die \emph{Lastimpedanz} am Knoten $k$ bezeichnet, vereinfacht sich \eqref{eq:eta} zu~\cite{Salazar2024,Giraldo2022}
\begin{equation}
	\label{eq:eta_physical}
	\eta = \left\|\mathbf{Z}_B\operatorname{diag}(\mathbf{1} \oslash \mathbf{z}_L)\right\|_1 \approx \max_{k \in \Omega_d}\frac{|Z_{\text{Th},k}|}{|z_{L,k}|}
\end{equation}
mit der Thévenin-Impedanz $Z_{\text{Th},k} = [\mathbf{Z}_B]_{kk}$ aus \eqref{eq:z_thev} reduziert werden.
Wenn die Hinreichende Bedingung erfüllt wird, gilt 
\begin{equation}
	\label{eq:physical_condition}
	|Z_{\text{Th},k}| < |z_{L,k}|
\end{equation}
für alle Lastknoten in Netz.
Physikalisch bedeutet dies, dass Lastwiderstand größer sein muss, als die Thévenin-Impedanz (Innenwiderstand). 
Diese Charakteristik ist genau im hochspannungsseitigen Betriebspunkt des elektrischen Netzes erfüllt \cite{Giraldo2022}.

%Die FPI konvergiert genau dann, wenn an jedem Lastknoten die Thévenin-Impedanz kleiner als die Lastimpedanz ist. Dies ist der praktische, hochspannungsseitige Betriebspunkt. 

%\paragraph{Zusammenhang mit der Existenzbedingung.}
%Die Kontraktionsbedingung \eqref{eq:physical_condition} fällt im Grenzfall zusammen mit der in Abschnitt~\ref{subsec:existenz} eingeführten Existenzbedingung \eqref{eq:feasibility} von Bolognani und Zampieri~\cite{bolognani2016}: Existiert eine praktische Lastflusslösung, so ist die FPI eine Kontraktion und konvergiert; existiert keine Lösung, so kann auch keine Kontraktion vorliegen. Damit ist die FPI in doppelter Hinsicht robust: Sie konvergiert genau dann, wenn eine physikalische Lösung existiert, und sie konvergiert stets zum \emph{richtigen} (hochspannungsseitigen) Betriebspunkt. \todo{macht gerade keinen sinn wenn es auch für große Eta konvergiert :'/}

%\subsection{Interpretation und Vergleich mit Newton-Raphson}
%\label{subsec:fpi_interpretation}
%
%
%Im Gegensatz zum NR-Verfahren mit quadratischer Konvergenz weist die FPI nur eine \emph{lineare} Konvergenzrate auf. Für typische Verteilnetze mit moderatem Belastungsgrad liegt $\eta$ im Bereich $0{,}05$--$0{,}30$~\cite{Giraldo2022}, was zu 5--15 Iterationen bei Toleranz $\varepsilon = 10^{-8}$ führt. Der scheinbare Nachteil höherer Iterationszahl wird jedoch durch den drastisch reduzierten Aufwand pro Iteration überkompensiert (siehe Kapitel~\ref{sec:tpf}).
%
%\paragraph{Robustheit gegenüber Startwert.}
%Ein zentraler Vorteil der FPI liegt in ihrer \emph{globalen} Konvergenz: Nach dem Banachschen Fixpunktsatz konvergiert die Iteration für \emph{jeden} Startwert $\mathbf{v}_d^{(0)}$ zur einzigen Lösung, sofern $\eta < 1$. Numerische Untersuchungen in \cite{salazar2024} bestätigen dies: Während NR je nach Startwert in den nicht-operativen Niederspannungsast konvergieren kann, findet die FPI stets den hochspannungsseitigen Betriebspunkt. Dies eliminiert die praktisch relevante Notwendigkeit eines guten Startpunktes (\emph{warm start}), was insbesondere bei massenparallelen Szenarienrechnungen mit stark variierenden Lastprofilen vorteilhaft ist.
%
%\paragraph{Vergleichende Bewertung.}
%Tabelle~\ref{tab:nr_fpi} fasst die wesentlichen Unterschiede zwischen NR und FPI zusammen.
%
%\begin{table}[h]
%	\centering
%	\caption{Gegenüberstellung Newton-Raphson und Fixpunkt-Iteration}
%	\label{tab:nr_fpi}
%	\begin{tabular}{lll}
%		\toprule
%		Eigenschaft & Newton-Raphson & Fixpunkt-Iteration \\
%		\midrule
%		Konvergenzordnung & quadratisch & linear \\
%		Typische Iterationen & 3--5 & 5--15 \\
%		Systemmatrix & $\mathbf{J}(\mathbf{x}^{(n)}, \mathbf{s})$ variabel & $\mathbf{Z}_B$ konstant \\
%		Faktorisierung & pro Iteration & einmalig \\
%		Startwertabhängigkeit & lokal konvergent & global konvergent \\
%		Konvergenzgarantie & keine & bei $\eta < 1$ \\
%		Konvergenzast & Start-abhängig & stets Hochspannung \\
%		Parallelisierung Szenarien & schwierig & natürlich (BLAS) \\
%		PV-Integration & natürlich (Jacobi) & problematisch \\
%		\bottomrule
%	\end{tabular}
%\end{table}
%
%Der letzte Punkt -- die problematische PV-Integration -- ist der methodische Kern dieser Arbeit und wird in Kapitel~\ref{sec:pv_methoden} ausführlich behandelt: Die Konstanz von $\mathbf{Z}_B$, die den Schlüsselvorteil der FPI ausmacht, ist gerade die Eigenschaft, die durch die naive Einführung von PV-Knoten (mit variablem $Q$) zerstört wird.
%
%\subsection{Konsequenzen für die Tensorformulierung}
%\label{subsec:fpi_eigenschaften}
%
%Für die im folgenden Kapitel eingeführte Tensorformulierung sind aus der FPI-Herleitung drei Eigenschaften von zentraler Bedeutung:
%
%\begin{enumerate}
%	\item \textbf{Konstante Systemmatrix.} Die Matrix $\mathbf{Z}_B = \mathbf{Y}_{dd}^{-1}$ hängt weder vom Iterationszustand $\mathbf{v}_d^{(n)}$ noch vom Lastszenario $\mathbf{s}_d$ ab. Sie kann einmalig vorab berechnet (dense) oder LU-faktorisiert (sparse) werden und für alle $\tau$ Szenarien und alle Iterationen wiederverwendet werden.
%	
%	\item \textbf{Einheitliche Iterationsstruktur.} Die Iteration \eqref{eq:fpi_zbus} besteht aus einer Matrix-Vektor-Multiplikation und einer elementweisen Operation. Beide lassen sich für $\tau$ Szenarien parallel als Matrix-Matrix-Multiplikation ausdrücken:
%	\begin{equation}
%		\label{eq:fpi_tensor}
%		\dot{\mathbf{V}}^{(n+1)} = \mathbf{Z}_B\left(\dot{\mathbf{S}}^{*} \oslash \dot{\mathbf{V}}^{*\,(n)}\right) + \mathbf{W}
%	\end{equation}
%	mit $\dot{\mathbf{V}}, \dot{\mathbf{S}} \in \mathbb{C}^{b\varphi \times \tau}$. Diese Struktur ist ideal für BLAS-basierte Bibliotheken (\texttt{OpenBLAS}, \texttt{Intel MKL}) sowie für GPU-Beschleunigung mittels CUDA/cuBLAS.
%	
%	\item \textbf{Konvergenzindikator $\eta$.} Über \eqref{eq:eta} bzw.\ \eqref{eq:eta_physical} lässt sich \emph{vor} Beginn der Iteration prüfen, ob und wie schnell die FPI konvergieren wird. Dies erlaubt die Validierung der Feasibility eines Szenarios ohne tatsächliche Iteration und dient in dieser Arbeit als zentrale Kennzahl zur Bewertung der Robustheit (vgl.\ Kap.~\ref{sec:ergebnisse}).
%\end{enumerate}
%
%\bigskip
%\noindent
%Zusammenfassend wurde die Fixpunkt-Iteration als explizite Iterationsvorschrift mit konstanter Systemmatrix $\mathbf{Z}_B$ hergeleitet. Über den Banachschen Fixpunktsatz wurde bewiesen, dass die Iteration bei $\eta < 1$ global gegen die eindeutige, physikalisch sinnvolle Lösung konvergiert. Die Kontraktionsbedingung wurde physikalisch als Thévenin-/Lastimpedanz-Verhältnis interpretiert und mit der Existenzbedingung des Lastflussproblems in Verbindung gebracht. Die drei zentralen Eigenschaften -- Konstanz der Systemmatrix, BLAS-fähige Iterationsstruktur und quantitative Konvergenzvorhersage -- bilden die Grundlage für die im folgenden Kapitel dargestellte Tensor Power Flow-Formulierung.



